
\section{Introduction}
Temporal graph learning represents a graph as an ordered stream of event interactions. It is widely used for edges that appear over time, such as paper citations, co-authorships, patent assignments, and software dependency events. 
Recent models such as TGAT, TGN, DyGFormer, GraphMixer, and RepeatMixer have shown that temporal attention, event memory, neighbor co-occurrence, and repeated interaction patterns are important for link prediction in continuous-time graphs \cite{xu2020tgat,rossi2020tgn,yu2023dygformer,cong2023graphmixer,zou2024repeatmixer}. 
\rev{In parallel, heterogeneous graph learning has shown that node and relation types are essential when similar topologies carry different semantics \cite{wang2019heterogeneous,hu2020hgt}.}

This paper studies temporal link prediction on dynamic heterogeneous and text-attributed graphs. 
This setting is significant for academic data mining and technology intelligence.
In scholarly networks, papers, authors, venues, and topics form typed relations that evolve with time. 
In patent networks, assignees, inventors, patents, and IPC/CPC codes interact through multiple technical ownership relations. 
In software ecosystems, users, repositories, packages, and issues are connected by typed temporal events. 
\rev{These applications require models that jointly use textual attributes, time-dependent structural context, and complex typed relations.}
\begin{figure}[t]
\centering
\includegraphics[width=\linewidth]{figures/1-existmethod.png}
\caption{Illustration of existing methods and gaps for temporal link prediction.
}
\label{fig:intro_results}
\vspace{-2mm}
\end{figure}

The task is challenging for three reasons. 
First, relation types have different meanings and predictive roles. A citation edge, a coauthor edge, and an assignee--technology edge should not share the same message function. 
Second, relation behavior is non-stationary. A relation that is informative in one period may become less informative later. 
Third, future links depend on recent temporal neighborhoods rather than isolated endpoints, especially when target edges are sparse or partially unseen.
However, existing approaches address these issues partially. Methods for homogeneous temporal graphs model interaction sequences, but usually ignore typed entities and multi-relational structures. 
Methods for heterogeneous or text-attributed graphs capture node types and attributes, but relation types are often used as fixed labels or static embeddings. 
As shown in \ref{fig:intro_results}, this leaves relation evolution and relation-aware local temporal context unexplored. 

To address these limitations, we propose DHC, a dynamic heterogeneous context-learning framework. 
DHC maintains node memory and relation memory, generates relation-conditioned temporal messages, 
encodes recent local neighborhoods, and fuses them with a sequential event-memory branch. 
The model follows a unified event-stream protocol. Context edges update dynamic states, while target edges provide supervised labels. 
This is useful for academic and patent graphs because auxiliary relations, such as authorship, venue, topic, and technology-paper links, can provide evidence for predicting a
specific target relation. 

The main contributions are summarized as follows:
\begin{itemize}[leftmargin=*]
\item We introduce a temporal heterogeneous link prediction framework that jointly models text-derived node attributes, relation-aware messages, node-relation dual memory, and recent-neighbor temporal context.
\item \rev{We design relation-time message generation and dynamic relation states, allowing relation semantics to influence both state evolution and link scoring.}
\item \rev{Experiments on four processed temporal heterogeneous datasets from scholarly, patent, and software domains evaluate DHC in both transductive and inductive settings.}
\end{itemize}
